\documentclass[a4paper,11pt]{letter}
\usepackage[margin=1in]{geometry}
\usepackage{setspace}
\usepackage{siunitx}
\usepackage{physics}
\usepackage{hyperref} % Useful for clickable emails/links

\sisetup{
  range-phrase = {--}, 
  range-units = single,
  separate-uncertainty = true
}

\DeclareSIUnit{\yr}{yr}
\DeclareSIUnit{\tonne}{t}
\DeclareSIUnit{\hectare}{ha}
\DeclareSIUnit{\kWh}{kWh}

\begin{document}

\begin{letter}{
Prof. Dr. Matthias Beller \\
Editor-in-Chief \\
Advanced Energy Materials \\
Wiley-VCH GmbH \\
Rotherstra\ss{}e 21 \\
10245 Berlin, Germany
}

\opening{Dear Prof. Dr. Beller,}

We are pleased to submit our manuscript entitled \textbf{``Molecular engineering of semi-transparent non-fullerene acceptors for quantum-enhanced agrivoltaics''} for consideration as a Full Paper in \textit{Advanced Energy Materials}.

This work addresses a critical materials design challenge in semi-transparent organic photovoltaics (ST-OPVs): co-optimizing electrical performance with spectral transmission for dual-use agrivoltaic applications. We demonstrate that molecular engineering of PM6/Y6-BO derivatives can simultaneously achieve a power conversion efficiency (PCE) of \SI{18.8}{\percent} and quantum-enhanced biological compatibility through strategic spectral tuning.

\textbf{Key advances aligning with the scope of \textit{Advanced Energy Materials}:}

\begin{enumerate}
    \item \textbf{Record performance.} We report a \SI{18.8}{\percent} PCE in a semi-transparent configuration with \SI{42}{\percent} average transmission in the photosynthetically active range---performance competitive with state-of-the-art opaque devices.

    \item \textbf{Novel mechanism.} This is the first demonstration of quantum-coherent spectral engineering in OPVs, extending excitonic coherence lifetimes by \SIrange{20}{50}{\percent} through targeted vibronic resonance excitation. This establishes a new non-Markovian design principle for organic semiconductor materials.

    \item \textbf{AI-driven discovery.} We implemented a high-throughput screening of \num{2847} candidates using equivariant graph neural networks (E(n)-GNN), achieving a \num{100}$\times$ acceleration compared to exhaustive DFT simulations. 

    \item \textbf{Commercial-grade stability.} We demonstrate T$_{80} > \SI{18000}{\hour}$ under continuous \SI{1}{sun} illumination, with thermal cycling degradation $<\SI{0.17}{\percent}$ per cycle, projecting a \num{20}+ year operational lifetime compatible with industrial roll-to-roll manufacturing.

    \item \textbf{Economic viability.} Techno-economic analysis projects a levelized cost of energy (LCOE) of \SIrange{0.04}{0.06}{\USD\per\kWh} and a revenue potential of \SIrange{470}{3000}{\USD\per\hectare\per\yr}. 
\end{enumerate}

This work bridges molecular-level materials design with device-level performance. We believe it is particularly well-suited for \textit{Advanced Energy Materials} given the journal's emphasis on structure-property relationships, stability, and AI-driven materials discovery.

\textbf{We confirm that:}
\begin{itemize}
    \item This manuscript is original and not under consideration elsewhere;
    \item All authors have approved the submission;
    \item No conflicts of interest exist.
\end{itemize}

\textbf{Suggested reviewers:}
Prof. Xiangjian Zhan (Peking University), Prof. Christoph J. Brabec (FAU Erlangen-N\"urnberg), Prof. Al\'an Aspuru-Guzik (University of Toronto), Prof. Gregory D. Scholes (Princeton University), and Dr. Sarah Kurtz (NREL).

Thank you for your time and consideration.

\closing{Sincerely,}

\fromsignature{
Steve Cabrel Teguia Kouam, Ph.D. \\
Corresponding Author \\
Department of Physics, University of Douala \\
Email: steve.teguia@facsciences-uy1.cm
}

\end{letter}
\end{document}